\SourceLeaf{345}
\par\medskip\Needspace{4\baselineskip}\pdfbookmark[2]{Art. 7. 8.}{nachlass-biquadratic-III-art-7-8}\centerline{\textbf{7. 8.}}\nobreak\smallskip
This rule is general, whatever number $M$ denotes. For the case that will be the immediate subject of our investigation, where $M$ is assumed to be odd and of the form $1+(2+2i)N$, a slightly modified rule is more convenient.

Suppose the numbers $f$ are again divided into $4$ classes. In the first we place the $(h)$, whose doubles are also found in $f$; in the second, $h'$, we count those whose doubles $2h'$ belong to $f'$; similarly, $h''$ and $h'''$ denote those whose doubles belong to $f''$ and $f'''$. Thus the decident $\varepsilon$ is
\[
\varepsilon=\Sigma\Theta\frac{2hM}{m}+\Sigma\Theta\frac{2h'M}{m}+\Sigma\Theta\frac{2h''M}{m}+\Sigma\Theta\frac{2h'''M}{m}
\]
We call the collection of all $2h$ and $-2h''+(1+i)m$ $H$

and the collection of all $-i(2h'-m)$ and $i(2h'''-im)$ $H'$

Thus $H$ and $H'$ comprise all $f$; the former are the even ones, the latter the odd ones.

The following relations must also be applied
\[
\begin{aligned}
\Theta iN&=1+\Theta N\\
\Theta(-N)&=2+\Theta N\\
\Theta(-iN)&=3+\Theta N\\
\Theta(N+1)&=1-\Theta N\\
\Theta(N+1+i)&=2+\Theta N\\
\Theta(N+i)&=3-\Theta N
\end{aligned}
\]
consequently
\[
\begin{aligned}
\Theta\frac{(-2h'i+mi)M}{m}&=3-\Theta\frac{-2h'iM}{m}=-\Theta\frac{2h'M}{m}\\
\Theta\frac{(-2h''+m(1+i))M}{m}&=2+\Theta\frac{-2h''M}{m}=\Theta\frac{2h''M}{m}\\
\Theta\frac{(2h'''i+m)M}{m}&=1-\Theta\frac{2h'''iM}{m}=-\Theta\frac{2h'''M}{m}
\end{aligned}
\]
and
\[
\begin{aligned}
\varepsilon={}&\Sigma\Theta\frac{2hM}{m}-\Theta\frac{(-2h'i+mi)M}{m}\\
&+\Theta\frac{(-2h''+m(1+i))M}{m}-\Theta\frac{(2h'''i+m)M}{m}\\
={}&\Sigma\Theta\frac{HM}{m}-\Sigma\Theta\frac{H'M}{m}
\end{aligned}
\]

\SourceLeaf{346}
or
\[
\varepsilon=\Sigma\pm\Theta\frac{fM}{m}
\]
where the upper sign is to be taken for even $f$, the lower for odd.

\Article{9}
Let, in general, $f=\xi+\eta i$. The numbers $\xi,\eta$ are restricted within certain bounds by the condition that $f$ be a least residue of $m$, or that, on putting $\frac{f}{m}=x+yi$, $x$ and $y$ must lie between the bounds $0$ and $\frac12$. Definite rules for these bounds could easily be given by distinguishing the different cases. If we give $\eta$ a particular value, $\xi$ will in turn have its particular bounds. For example, if we assume that $a$ is negative and $b$ positive, then, since
\[
\begin{aligned}
x&=\frac{a\xi+b\eta}{aa+bb}\\
y&=\frac{a\eta-b\xi}{aa+bb}
\end{aligned}
\]
the following conditions must hold:
\[
\begin{array}{rll}
\mathrm{I.}&\text{for }x\text{ to be positive}&\xi<-\frac{b}{a}\eta\\[3pt]
\mathrm{II.}&\text{for }y\text{ to be positive}&\xi<\frac{a}{b}\eta\\[3pt]
\mathrm{III.}&\text{for }x<\frac12\text{ to hold}&\xi>\frac{aa+bb-2b\eta}{2a}\\[3pt]
\mathrm{IV.}&\text{for }y<\frac12\text{ to hold}&\xi>\frac{2a\eta-aa-bb}{2b}
\end{array}
\]

For positive $\eta$, the second condition already includes the first; for negative $\eta$, the reverse holds. Likewise, the third condition is already contained in the fourth,
\[
\begin{array}{ll}
\text{if}&\eta<\tfrac12(a+b)\\
\text{and conversely if}&\eta>\tfrac12(a+b)
\end{array}
\]

We do not, however, need to consider separately all eight cases that can occur here. For a particular value of $\eta$, we merely denote the smaller bound of $\xi$ by $\xi^0$, the larger by $\xi^{00}$, and remark only that at these boundary values one always has either $x=0,y=0,x=\frac12,y=\frac12$, and specifically that the following must hold:
\SourceLeaf{347}
\[
\begin{array}{l|l|l}
\text{if}&\text{at the \emph{upper} bound}&\text{at the \emph{lower} bound}\\
a\text{ pos. }b\text{ positive}&x=\frac12\text{ or }y=0&x=0\text{ or }y=\frac12\\
a\text{ neg. }b\text{ positive}&x=0\text{ or }y=0&x=\frac12\text{ or }y=\frac12\\
a\text{ neg. }b\text{ negative}&x=0\text{ or }y=\frac12&x=\frac12\text{ or }y=0\\
a\text{ pos. }b\text{ negative}&x=\frac12\text{ or }y=\frac12&x=0\text{ or }y=0
\end{array}
\]
For brevity we shall distinguish these four cases by saying that, in the first, $m$ belongs to the first quadrant, in the second to the second, and so on.

\Article{10}
We now wish to examine more closely the sum of all $\pm\Theta\frac{fM}{m}$ for which $\eta$ has a particular value. As $\xi$ gradually increases continuously from its smallest value $\xi^0$ to its largest $\xi^{00}$,
\[
\frac{(\xi+\eta i)M}{m}=X+Yi
\]
will also change according to the law of continuity. If $\frac{M}{m}$ lies in the first quadrant, both $X$ and $Y$ will keep increasing; if $\frac{M}{m}$ lies in the second quadrant, $X$ will keep decreasing and $Y$ increasing. In the third quadrant the reverse of the first case occurs, and in the fourth the reverse of the second. But the integer contained in $X+iY$ will change in jumps, as either $[X]$ or $[Y]$ changes by one unit. Let the values of $\xi$ at which such a transition occurs, that is, at which either $X$ or $Y$ becomes an integer, be, in order,
\[
\xi',\ \xi'',\ \xi'''\dotsfour\xi^n
\]

Here it must be remarked that neither these values nor $\xi^0$ and $\xi^{00}$ can be integers, except for $\eta=0$, when either $\xi^0$ or $\xi^{00}=0$. Now let
\[
\begin{aligned}
\Theta\frac{(\xi'+\eta i)M}{m}-\Theta\frac{(\xi^0+\eta i)M}{m}&=\delta'\quad\text{(to be expressed differently)}\\
\Theta\frac{(\xi''+\eta i)M}{m}-\Theta\frac{(\xi'+\eta i)M}{m}&=\delta''\\
&\text{etc.}\\
\Theta\frac{(\xi^{00}+\eta i)M}{m}-\Theta\frac{(\xi^{n-1}+\eta i)M}{m}&=\delta^n
\end{aligned}
\]
\SourceLeaf{348}
Then it is easy to see, since between $\xi^0$ and $\xi'$ there lie $[\frac12\xi']-[\frac12\xi^0]$ even integers and $[\frac12\xi'+\frac12]-[\frac12\xi^0+\frac12]$ odd integers, and so on, that, considering only the particular value of $\eta$,
\[
\begin{aligned}
(\pm1)\Sigma\Theta\frac{fM}{m}
={}&\{[\tfrac12\xi']-[\tfrac12\xi^0]-[\tfrac12\xi'+\tfrac12]+[\tfrac12\xi^0+\tfrac12]\}\\
&\qquad\pdot\Theta\frac{(\xi^0+\eta i)M}{m}\\
&+\{[\tfrac12\xi'']-[\tfrac12\xi']-[\tfrac12\xi''+\tfrac12]+[\tfrac12\xi'+\tfrac12]\}\\
&\qquad\pdot\{\Theta\frac{(\xi^0+\eta i)M}{m}+\delta'\}\\
&+\{[\tfrac12\xi''']-[\tfrac12\xi'']-[\tfrac12\xi'''+\tfrac12]+[\tfrac12\xi''+\tfrac12]\}\\
&\qquad\pdot\{\Theta\frac{(\xi^0+\eta i)M}{m}+\delta'+\delta''\}\\
&+\text{etc.}\\
&+\{[\tfrac12\xi^{00}]-[\tfrac12\xi^n]-[\tfrac12\xi^{00}+\tfrac12]+[\tfrac12\xi^n+\tfrac12]\}\\
&\qquad\pdot\{\Theta\frac{(\xi^0+\eta i)M}{m}+\delta'+\delta''+\dotstwo+\delta^n\}
\end{aligned}
\]
\[
\begin{aligned}
={}&-([\tfrac12\xi^0]-[\tfrac12\xi^0+\tfrac12])\pdot\Theta\frac{(\xi^0+\eta i)M}{m}\\
&-([\tfrac12\xi']-[\tfrac12\xi'+\tfrac12])\pdot\delta'\\
&-([\tfrac12\xi'']-[\tfrac12\xi''+\tfrac12])\pdot\delta''\\
&-\text{etc.}\\
&-([\tfrac12\xi^n]-[\tfrac12\xi^n+\tfrac12])\pdot\delta^n\\
&+([\tfrac12\xi^{00}]-[\tfrac12\xi^{00}+\tfrac12])\pdot\Theta\frac{(\xi^{00}+\eta i)M}{m}
\end{aligned}
\]
(where the upper sign applies for even $\eta$, the lower for odd.)

\Needspace{17\baselineskip}The numbers $\delta',\delta'',\delta'''$, and so on, can have no values other than $+1$ and $-1$. The value $+1$ is obtained by $\delta'$ when, denoting the values of $X,Y$ belonging to $\xi'$ by $X',Y'$,
\[
\begin{array}{l|l}
\frac{M}{m}\text{ in quadrant 1.}&\frac{M}{m}\text{ in quadrant 2.}\\
X'\text{ even integer}&X'\text{ even integer}\\
\text{and }[Y]\text{ odd}&[Y]\text{ even}\\\hline
Y'\text{ even integer}&Y'\text{ even integer}\\
\text{and }[X]\text{ even}&[X]\text{ even}
\end{array}
\]
\[
\begin{array}{l|l}
\frac{M}{m}\text{ in quadrant 3.}&\frac{M}{m}\text{ in quadrant 4.}\\
X'\text{ even integer}&X'\text{ even integer}\\
\text{and }[Y]\text{ even}&[Y]\text{ odd}\\\hline
Y'\text{ even integer}&Y'\text{ even integer}\\
\text{and }[X]\text{ odd}&[X]\text{ odd}
\end{array}
\]

\SourceLeaf{349}
Whenever one of these conditions changes to its opposite, $\delta=-1$; whenever both change, $\delta=+1$ remains.

\Article{11}
The following formulas provide a more convenient overview of these calculations:
\[
\begin{aligned}
\text{here }m&=a+bi,\qquad aa+bb=d\\
M&=A+Bi,\qquad AA+BB=D\\
\frac{dM}{m}&=\alpha+\beta i,\qquad\alpha=aA+bB,\quad\beta=aB-bA\\
\frac{\xi+i\eta}{m}&=x+iy,\qquad M(x+iy)=X+iY
\end{aligned}
\]

If $\eta$ and $X$ are given, then
\[
\begin{array}{rl}
1.&\xi=\frac{\beta\eta}{\alpha}+\frac{dX}{\alpha}\\[3pt]
2.&Y=\frac{D\eta}{\alpha}+\frac{\beta X}{\alpha}
\end{array}
\]
If $\eta$ and $Y$ are given, then
\[
\begin{array}{rl}
3.&\xi=-\frac{\alpha\eta}{\beta}+\frac{dY}{\beta}\\[3pt]
4.&X=-\frac{D\eta}{\beta}+\frac{\alpha Y}{\beta}
\end{array}
\]
If $\eta$ and $x$ are given, then
\[
\begin{array}{rl}
5.&\xi=-\frac{b\eta}{a}+\frac{dx}{a}\\[3pt]
6.&X=-\frac{B\eta}{a}+\frac{\alpha x}{a}\\[3pt]
7.&Y=\frac{A\eta}{a}+\frac{\beta x}{a}
\end{array}
\]
If $\eta$ and $y$ are given, then
\[
\begin{array}{rl}
8.&\xi=\frac{a\eta}{b}-\frac{dy}{b}\\[3pt]
9.&X=\frac{A\eta}{b}-\frac{\alpha y}{b}\\[3pt]
10.&Y=\frac{B\eta}{b}-\frac{\beta y}{b}
\end{array}
\]

\Article{12}
The rule of Article $10$ can now be expressed as follows. Giving $\eta$ a particular value, we have
\[
\Sigma\pm\Theta\frac{fM}{m}=k^0\Theta(X^0+Y^0i)-k^{00}\Theta(X^{00}+Y^{00}i)+\Sigma k
\]

\SourceLeaf{350}
Here $k^0=0$ when $[\xi^0]$ is even; $k^0=+1$ when $[\xi^0]$ is odd and $\eta$ even; $k^0=-1$ when $[\xi^0]$ is odd and $\eta$ odd. $k^{00}$ is determined in the same way by $[\xi^{00}]$ and $\eta$. Finally, $\Sigma k$ is a sum of as many numbers as there are, between $\xi=\xi^0$ and $\xi=\xi^{00}$, integer values of $X$ or $Y$; in each case $k=0$ when the corresponding $[\xi]$ is even, and $=\pm1$ when $[\xi]$ is odd. The sign is determined as follows. If $X$ is an integer, then $k=1$ when all the following hold:
\[
\begin{array}{ll}
\eta&\text{even}\\
X&\text{even}\\
\relax[Y]&\text{even}\\
\frac{M}{m}&\text{in the second or third quadrant, i.e. }\alpha\text{ negative}
\end{array}
\]

If one or three of these conditions are absent, then $k=-1$; if two or all four are absent, $k=1$ remains. If, on the other hand, $Y$ is an integer, then $k=1$ when all, two, or none of the $4$ conditions
\[
\begin{array}{ll}
\eta&\text{even}\\
Y&\text{even}\\
\relax[X]&\text{even}\\
\frac{M}{m}&\text{in the first or second quadrant, i.e. }\beta\text{ positive}
\end{array}
\]
are satisfied.

\Article{13}
We must now consider separately the cases in which $\xi^0$ or $\xi^{00}$ (or $X^0,Y^0,X^{00},Y^{00}$) is an integer. There are four cases to distinguish here, taking $a$ and $A$ to be odd.

I. If $m$ lies in the first quadrant, then for $x=\frac12,y=0$, $\eta=\frac12b$ is an integer. Then $Y^{00}=\frac12B$ is an integer, and $\Theta(X^{00}+Y^{00}i)$ will be $=\Theta(\frac12A+\frac12Bi)$ only when $\beta$ is negative. When $\beta$ is positive, $\Theta(\frac12A+(\frac12B-1)i)$ is to be taken instead.

II. If $m$ lies in the second quadrant, then for $x=0,y=0$ one has $\eta=0$. For this value of $\eta$, $X^{00}=0,Y^{00}=0$. This gives

\SourceLeaf{351}
\[
\begin{array}{rll}
\Theta(X^{00}+Y^{00}i)=&2,&\text{according as }\dfrac{M}{m}\text{ lies in quadrant }1.\\
&3,&\hfill2.\\
&0&\hfill3.\\
&1&\hfill4.\text{, and }k^{00}=1
\end{array}
\]

III. If $m$ lies in the third quadrant, then for $x=\frac12,y=0$, $\eta=\frac12b$ is an integer, for which $X^0+Y^0i=\frac12A+\frac12Bi$. Then put
\[
\Theta(X^0+Y^0i)=\Theta(\tfrac12A+(\tfrac12B-1)i)
\]
whenever $\beta$ is negative.

IV. If $m$ lies in the fourth quadrant, then for $\eta=0$ put
$\Theta(X^0+Y^0i=0,1,2,3$, according as $\frac{M}{m}$ lies in quadrant $1.\ 2.\ 3.\ 4.$;
$k^0=0$.

\Article{14}
The preceding investigations now yield the following determination of the decident.

Collect all values of $x$ and $y$ that lie \emph{within} the bounds $0$ and $\frac12$ and for which either $\eta$ and $X$ or $\eta$ and $Y$ are integers. For each $x+iy$, use the rules of Article $12$ to determine the value of $k$.

Further, collect all values on the boundaries, that is, where either $x=0$ or $\frac12$ while $y$ lies \emph{between} $0$ and $\frac12$, or $y=0$ or $=\frac12$ while $x$ lies \emph{between} $0$ and $\frac12$, such that $\eta$ is an integer and $[\xi]$ odd. Determine the corresponding $l$ as follows. Let $\Theta M(x+yi)=\pm\theta$, with the upper sign for even, the lower for odd $\eta$
\[
\begin{array}{c|c|c|c|c}
&\multicolumn{4}{c}{\text{then for }m\text{ in}}\\
\text{for}&\text{quadrant }1.&\text{quadrant }2.&\text{quadrant }3.&\text{quadrant }4.\\
y=0&l=-\theta&l=-\theta&l=+\theta&l=+\theta\\
x=\frac12&l=-\theta&l=+\theta&l=+\theta&l=-\theta\\
y=\frac12&l=+\theta&l=+\theta&l=-\theta&l=-\theta\\
x=0&l=+\theta&l=-\theta&l=-\theta&l=+\theta
\end{array}
\]

\SourceLeaf{352}
More briefly,
\[
l=\pm\theta,
\]
\[
\begin{array}{ll}
\text{the sign is the same as that of }a&\text{if }x=0\\
\text{the opposite}&\text{if }x=\frac12\\
\text{the same as that of }b&\text{if }y=\frac12\\
\text{the opposite}&\text{if }y=0
\end{array}
\]

\Needspace{22\baselineskip}To $\Sigma k+\Sigma l$ one must then add
\begingroup\setlength{\arraycolsep}{1pt}
\[
\left.\begin{array}{ll}
\text{if }m\text{ lies in the second quadrant:}&2,1,0,3\\
\text{if }m\text{ lies in the fourth quadrant:}&0,
\end{array}\right\}
\begin{array}{l}\text{according as }\frac{M}{m}\text{ lies in quadrant}\\1.\ 2.\ 3.\ 4.\text{}\end{array}
\]\endgroup
if $m$ lies in the first quadrant and $\frac12(a-1)$ is odd,
\[
\begin{array}{ll}
\Theta(\frac12A+(\frac12B-1)i)&\text{if }\beta\text{ positive}\\
\Theta(\frac12A+\frac12Bi)&\text{if }\beta\text{ negative}
\end{array}
\]
if $m$ lies in the third quadrant and $\frac12(a-1)$ is odd,
\[
\begin{array}{ll}
-\Theta(\frac12A+\frac12Bi)&\text{if }\beta\text{ positive}\\
-\Theta(\frac12A+(\frac12B-1)i)&\text{if }\beta\text{ negative.}
\end{array}
\]
\BookRule{.3\textwidth}
